\section{Future Work}

Every optimization discovered by \tool{} is specific: it is not
parameterized by bitwidths, values of constants, or choice or ordering
of instructions.
%
On the other hand, optimizations implemented by compiler developers are
almost always generic along one or more of these dimensions.
%
It would be useful for \tool{} to automate some of this
generalization, for example emitting an Alive~\cite{Lopes15} pattern
for each optimization it discovers.
%
Generalization would improve our rankings by preventing important
optimizations that happen to have many different forms from hiding low
in the rankings where they are unlikely to be looked at.
%
Generalization would also present compiler developers with
optimizations that are more like the ones that they will presumably
implement, saving them the effort of generalization by hand.
%
In particular, it can be very difficult to write a sound and precise
precondition for an optimization, especially in the presence
of undefined behavior.


A middle-end superoptimizer cannot exploit
target-specific code sequences.
%
Our hypothesis is that future compilers should contain two
superoptimizers.
%
First, one in the middle-end that, like \tool{}, tries to generate
constants and perform other obviously-profitable transformations that
can be exploited by other middle-end passes.
%
Second, a target-aware backend superoptimizer that interacts with
instruction selection, scheduling, and register allocation.


Finally, we believe that \tool{} IR can be extracted from CompCert RTL
in the same way that we extract it from LLVM IR and Microsoft Visual
C++ IR\@.
%
Synthesizing optimizations is then trivial.
%
The research problem is to take the resulting equivalence proofs
emitted by a proof-producing solver and integrate them into CompCert's
overall proof.
